\documentclass[a4paper,10pt,fleqn]{article}
\usepackage[margin=1in]{geometry}
\usepackage{amssymb}
\usepackage{adjustbox}
\usepackage{natbib}
\usepackage[colorlinks=true,linkcolor=blue,citecolor=blue,urlcolor=blue]{hyperref}
\usepackage{fuzz}
\usepackage{schemapk}
\usepackage{zed-maths}
\usepackage{zed-proof}
\newdimen\savedleftskip
\title{Multi-Declaration Z Constructs --- Demo}
\author{txt2tex}
\hypersetup{pdftitle={Multi-Declaration Z Constructs --- Demo},pdfauthor={txt2tex},pdfcreator={txt2tex v1.6.3}}
\begin{document}

\maketitle

\section*{Background}

\noindent The constructs in this file all share a single grammatical shape: a Z paragraph that introduces two or more typed variables in one declaration list, with a conjunction predicate over their projections, optionally followed by a characteristic expression. Until the multi-decl parser fix, none of them parsed natively — users wrote them as raw LaTeX. They now compile to fuzz-clean output.

\bigskip

\begin{zed}[PlayerName, TournamentId, MatchId, VenueName]\end{zed}

\medskip

\setbox0=\vbox{%
\begin{schema}{Tournament}
tournamentId : TournamentId \\
venue : VenueName \\
tier : \nat
\where
tier \leq 5
\end{schema}%
}
\begin{schemapk}{Tournament}
\underline{tournamentId} : TournamentId \\
venue : VenueName \\
tier : \nat
\where
tier \leq 5
\end{schemapk}

\setbox0=\vbox{%
\begin{schema}{Player}
name : PlayerName \\
ranking : \nat
\end{schema}%
}
\begin{schemapk}{Player}
\underline{name} : PlayerName \\
ranking : \nat
\end{schemapk}

\setbox0=\vbox{%
\begin{schema}{Match}
matchId : MatchId \\
tournament : TournamentId \\
winner : PlayerName \\
loser : PlayerName
\end{schema}%
}
\begin{schemapk}{Match}
\underline{matchId} : MatchId \\
\underline{tournament} : TournamentId \\
winner : PlayerName \\
loser : PlayerName
\end{schemapk}

\section*{1. WYSIWYG line breaks for algebra operators}

\noindent Long relational-algebra chains can be broken at any operator. Natural newlines work; explicit \ continuation also works. Display position wraps in array; the chain reads top-to-bottom instead of running off the page.

\bigskip

\noindent
$\displaystyle
\begin{array}{l}
\mathrm{Project}_{winner, venue, tier}(\mathrm{Join}(Tournament, \\
\quad (\mathrm{Project}_{tournament}(\mathrm{Restrict}_{winner = `Ali'}(Match)))[tournamentId/tournament]))
\end{array}$

\bigskip

\section*{2. Multi-decl comprehension with conjunction predicate}

\noindent This is the canonical multi-typed relational-calculus form for any query that crosses two or more relations. Before the fix it required a raw LATEX block; now it parses natively.

\bigskip

\setbox0=\vbox{%
\begin{zed}
zS_1 == \{~ m : Match ; t : Tournament ; p : Player | m.tournament = t.tournamentId \land \\
\quad m.winner = p.name \land \\
\quad t.venue = `Centre Court' @ \\ (m.winner, t.venue, t.tier) ~\}
\end{zed}%
}
\noindent
$\displaystyle
\begin{array}{l}
\begin{array}{l}\{~ m : Match ; t : Tournament ; p : Player | m.tournament = t.tournamentId \land \\
\quad m.winner = p.name \land \\
\quad t.venue = `Centre Court' @ \\ \lblot~winner == m.winner, venue == t.venue, tier == t.tier~\rblot  ~\}\end{array}
\end{array}$

\bigskip

\section*{3. Multi-decl quantifier with conjunction predicate}

\noindent `$\forall s : T; c : U @ predicate$. body` works the same way --- nested Quantifier nodes covering the chained declarations.

\bigskip

\begin{zed}
\forall m : Match; t : Tournament | m.tournament = t.tournamentId \land t.tier \leq 2 @ m.winner \neq m.loser
\end{zed}

\section*{4. Multi-decl quantifier without characteristic expression}

\noindent When the predicate alone suffices --- no `.` separator, no body expression --- the comprehension defaults the characteristic tuple to the signature (Z RM §3.9).

\bigskip

\setbox0=\vbox{%
\begin{zed}
zS_2 == \{~ m : Match ; t : Tournament | m.tournament = t.tournamentId \land t.tier \leq 2 ~\}
\end{zed}%
}
\noindent
$\{~ m : Match ; t : Tournament | m.tournament = t.tournamentId \land t.tier \leq 2 ~\}$


\section*{5. Multi-decl mu (definite description)}

\noindent `mu` parses with multi-decl signatures the same way. The current generator emits nested `(\mu ... $\mid$ (\mu ...))` --- semantically equivalent but not Spivey-canonical. The same special-case treatment lambda just received is a follow-up for mu. Single-decl mu (shown here) emits cleanly.

\bigskip

\begin{zed}
topFinal == (\mu t : Tournament | t.tier = 1 @ t.venue)
\end{zed}

\section*{6. Multi-decl lambda --- Spivey-canonical form}

\noindent Lambdas of more than one variable now emit as a single \lambda token with semicolon-joined SchemaText, matching Z RM §3.12 and the Spivey convention. Previously the generator emitted nested lambdas (one per declaration).

\bigskip

\begin{zed}
roster == (\lambda m : Match; t : Tournament | m.tournament = t.tournamentId @ m.winner)
\end{zed}

\section*{7. Paren-wrap fix}

\noindent Fuzz requires `(\lambda ... @ ...)` --- parentheses around every lambda expression, at any position. The generator previously only wrapped when the lambda appeared inside another expression (`parent is not None`). This meant abbreviation RHS, set literals, sequence displays, and top-level zed predicates emitted bare `\lambda`, which fuzz rejected with "Opening parenthesis expected at symbol \lambda".

\bigskip

\medskip

\noindent The fix: drop the `parent is not None` guard. Wrapping is now unconditional when `$use\_fuzz = True$`, matching how `\mu` has always been handled. All thirteen lambda-paren tests now pass without xfail markers. Section 6 above also now fuzz-typechecks cleanly.

\bigskip

\medskip

\noindent Additional positions confirmed working: lambda as abbreviation RHS (`f == lambda ...`), inside a set literal (`{lambda ...}`), inside a sequence display (`$\langle lambda ... \rangle$`), as a function argument (`f(lambda ...)`), as a tuple component (`(lambda ..., 0)`), and as a unary-operator argument (`dom (lambda ...)`).

\bigskip

\begin{zed}
f == (\lambda x : \nat @ x + 1)
\end{zed}

\end{document}